% Figure: kinetic-energy split of a body rolling without slipping down an
% incline. A cylinder, a sphere, and a thin ring release from the same height
% and reach the bottom at different speeds because a larger share of their
% energy goes into rotation. The bar chart shows the translational and
% rotational fractions of the total kinetic energy for each shape.
\begin{figure}[htbp]
\centering
\begin{tikzpicture}[
  axis/.style={draw=engblack, -{Stealth}, thick},
  transbar/.style={draw=engblue, fill=engblue!25},
  rotbar/.style={draw=engred, fill=engred!25},
  label/.style={font=\small},
]
% y axis: fraction of kinetic energy 0..1, plotted 0..6 (1 unit = 1/6)
\draw[axis] (0,0) -- (0,6.6) node[anchor=south, label] {share of kinetic energy};
\draw[axis] (0,0) -- (9.2,0) node[anchor=west, label] {rolling body};
\foreach \yp/\yv in {0/0.0, 3/0.5, 6/1.0} {
  \draw (-0.08,\yp) -- (0.08,\yp);
  \node[font=\scriptsize, left] at (-0.12,\yp) {$\yv$};
}
% Three bodies. Translational fraction = 1/(1+beta), rotational = beta/(1+beta)
%   thin ring   beta = 1     -> trans 1/2 = 0.500, rot 0.500
%   cylinder    beta = 1/2   -> trans 2/3 = 0.667, rot 0.333
%   sphere      beta = 2/5   -> trans 5/7 = 0.714, rot 0.286
% bar bottoms at 0, stack translational (blue) then rotational (red).
% plotted height = 6 * fraction.
% ring at x=1.2..2.4
\fill[transbar] (1.2,0) rectangle (2.4,3.000);
\fill[rotbar]   (1.2,3.000) rectangle (2.4,6.000);
\node[font=\scriptsize, below] at (1.8,-0.15) {ring};
\node[font=\scriptsize, below] at (1.8,-0.55) {$\beta=1$};
% cylinder at x=4.0..5.2
\fill[transbar] (4.0,0) rectangle (5.2,4.000);
\fill[rotbar]   (4.0,4.000) rectangle (5.2,6.000);
\node[font=\scriptsize, below] at (4.6,-0.15) {cylinder};
\node[font=\scriptsize, below] at (4.6,-0.55) {$\beta=\tfrac12$};
% sphere at x=6.8..8.0
\fill[transbar] (6.8,0) rectangle (8.0,4.286);
\fill[rotbar]   (6.8,4.286) rectangle (8.0,6.000);
\node[font=\scriptsize, below] at (7.4,-0.15) {sphere};
\node[font=\scriptsize, below] at (7.4,-0.55) {$\beta=\tfrac25$};
% legend
\fill[transbar] (5.6,5.6) rectangle (5.9,5.9);
\node[font=\scriptsize, anchor=west] at (5.95,5.75) {translation};
\fill[rotbar] (5.6,5.0) rectangle (5.9,5.3);
\node[font=\scriptsize, anchor=west] at (5.95,5.15) {rotation};
\end{tikzpicture}
\caption[Kinetic-energy split of a rolling body]{The kinetic energy of a body
rolling without slipping divides between translation of its centre and
rotation about it, in a ratio fixed by the shape factor $\beta = I_G/(mr^2)$.
A thin ring ($\beta=1$) puts half its energy into spin; a solid cylinder
($\beta=\tfrac12$) a third; a solid sphere ($\beta=\tfrac25$) less than
three-tenths. Released from the same height, the sphere reaches the bottom
fastest and the ring slowest, because the ring must spend the largest share
of the released gravitational energy on rotation it cannot convert to speed.
The energy method reads the finishing speed straight off the split without
resolving the friction force that supplies the rolling constraint.}
\label{fig:vol03ch05:rolling-energy-split}
\end{figure}
